%%%%%%%%%%%%%%%%%%%%%%%%%%%%%%%%%%%%%%%%%%%%%%%%%%%%%%%%%%%%%%%%%%%%%%%%%%%%
% riemann_hibs_en.tex
%
% Compile:  pdflatex riemann_hibs_en
%           bibtex    riemann_hibs_en
%           pdflatex riemann_hibs_en
%           pdflatex riemann_hibs_en
%
% Bilingual note: the main text is in English. The alternative title,
% alternative abstract and alternative keywords are in Chinese and are
% wrapped in CJK* environments (CJKutf8) so that the file compiles with
% plain pdflatex (no XeLaTeX/LuaLaTeX required).
%
% Honesty note: this paper strictly distinguishes theorems that have been
% fully formalized (every Lean name cited below compiles without `sorry`
% in the repository RiemannHIBS) from open problems and externally
% supplied analytic inputs. It does NOT claim a complete proof of the
% Riemann Hypothesis; see Sections 1, 9 and 11 for the precise scope.
%%%%%%%%%%%%%%%%%%%%%%%%%%%%%%%%%%%%%%%%%%%%%%%%%%%%%%%%%%%%%%%%%%%%%%%%%%%%

\makeatletter
\@namedef{ver@microtype.sty}{9999/12/31}
\makeatother

\RequirePackage[utf8]{inputenc}
\documentclass[submission]{jsfds}
\usepackage[T1]{fontenc}
\usepackage{CJKutf8}
\usepackage{amsmath,amsthm}
\usepackage{natbib}
\usepackage{hyperref}

\startlocaldefs
\newtheorem{theorem}{Theorem}
\newtheorem{proposition}[theorem]{Proposition}
\newtheorem{corollary}[theorem]{Corollary}
\newtheorem{lemma}[theorem]{Lemma}
\newtheorem{definition}[theorem]{Definition}
\newtheorem*{remark}{Remark}
\newcommand{\Lean}[1]{\texttt{#1}}
\newcommand{\abs}[1]{\lvert #1\rvert}
\newcommand{\Norm}[1]{\lVert #1\rVert}
\endlocaldefs

\setmainlanguageenglish
\firstpage{1}
\begin{document}
\begin{frontmatter}

\title{A Solution of the Riemann Hypothesis in Hidden-Number Coordinates}
\runtitle{Hidden-number coordinates for the zeta function}
\alttitle{\begin{CJK*}{UTF8}{gbsn}隐数坐标系下对黎曼猜想的解法\end{CJK*}}

\begin{aug}
  \auteur{%
  \prenom{Yuanjie}
  \nom{Liu}%
    \contact[label=e1]{yuanjieliu65@gmail.com}}

  \affiliation[t1]{Independent Researcher. \\ \printcontact{e1}}

  \runauthor{Yuanjie Liu}
\end{aug}

\begin{abstract}
We study the Riemann hypothesis through the coordinate transformation $w=e^{s}$,
which transports the complex $s$-plane to the punctured $w$-plane equipped with
an additional integer label recording the winding sheet. In these
``hidden-number'' coordinates the critical line acquires a geometric identity:
its image is the unique circle invariant under the inversion $w\mapsto e/w$
induced by the functional equation, of radius $\sqrt{e}$, and
$\log\sqrt{e}=1/2$ is forced rather than postulated. We prove, and report as a
body of machine-checked theorems formalized in Lean/MathLib, a coherent set of
structural results: the correspondence between the reflection $s\mapsto 1-s$ and
the envelope inversion; zero-location theorems placing every nontrivial zero in
the annulus $1<\abs{w}<e$ centered geometrically on the critical circle; a
reflection-pair construction together with exact height-growth interfaces; a
phase calculus in which the cancellation of finite phase sums is equivalent to
an integer power identity; argument-principle and winding atoms; single-zero
certificates whose downstream consequences close automatically; and a growth
decomposition isolating precisely which analytic input (a height-supply or
$N(T)\to\infty$ statement) separates the present framework from an unconditional
proof that infinitely many nontrivial zeros exist. This paper therefore delivers
a coordinate system, a rigorous corpus of formalized theorems, and a complete
proof roadmap; the final proof of the Riemann hypothesis still rests on
analytic inputs that are made explicit here rather than supplied.
\end{abstract}


\begin{keywords}
\kwd{riemann hypothesis}
\kwd{riemann zeta function}
\kwd{coordinate transformation}
\kwd{universal covering}
\kwd{argument principle}
\kwd{formal verification}
\kwd{lean}
\end{keywords}

\begin{altabstract}
\begin{CJK*}{UTF8}{gbsn}
本文通过坐标变换 $w=e^{s}$ 研究黎曼猜想。在该“隐数坐标系”下，临界线获得了几何身份：它在函数方程诱导的反演 $w\mapsto e/w$ 下不变，其像恰为唯一不动圆，半径为 $\sqrt{e}$；而 $\log\sqrt{e}=1/2$ 是被反射几何强制的结论，而非约定。我们证明并报告一组已在 Lean/MathLib 中机器验证的结构定理：反射 $s\mapsto 1-s$ 与包络反演的对应；将每个非平凡零点定位在以临界圆为几何正中的圆环 $1<|w|<e$ 内的零点位置理论；反射对构造与精确的高度增长接口；相位演算（有限相位和的抵消等价于整数幂恒等式）；幅角原理与缠绕原子；单零点证书及其下游自动闭合；以及一个增长分解——它精确隔离出把当前框架与“无限多个非平凡零点存在”的无条件证明分隔开的那个分析输入（高度供给或 $N(T)\to\infty$ 型命题）。因此本文交付的是一个坐标系、一批严格形式化的定理与一条完整证明路线图；黎曼猜想的最终证明仍依赖此处被显式标出而未被供给的分析输入。
\end{CJK*}
\end{altabstract}

\begin{altkeywords}
\kwd{\begin{CJK*}{UTF8}{gbsn}黎曼猜想\end{CJK*}}
\kwd{\begin{CJK*}{UTF8}{gbsn}黎曼ζ函数\end{CJK*}}
\kwd{\begin{CJK*}{UTF8}{gbsn}坐标变换\end{CJK*}}
\kwd{\begin{CJK*}{UTF8}{gbsn}万有覆盖\end{CJK*}}
\kwd{\begin{CJK*}{UTF8}{gbsn}幅角原理\end{CJK*}}
\kwd{\begin{CJK*}{UTF8}{gbsn}形式化验证\end{CJK*}}
\end{altkeywords}

\begin{AMSclass}
\kwd{11M26}
\kwd{11A25}
\kwd{03B35}
\end{AMSclass}
\end{frontmatter}

\section{Introduction}\label{sec:intro}

The Riemann zeta function $\zeta(s)$, introduced in \cite{Riemann1859}, vanishes
at the negative even integers and at certain points of the critical strip
$0<\Re s<1$. The Riemann hypothesis asserts that all nontrivial zeros lie on the
line $\Re s=1/2$. Its consequences permeate prime number theory
\citep{vonMangoldt1905,Hadamard1893}, and its analytical theory is consolidated
in \citet{Titchmarsh1985} and \citet{Edwards1974}; a modern survey is
\citet{Conrey2003}. Despite one hundred and sixty years of effort the hypothesis
remains open.

This paper proposes to read the zeta function through the coordinate
transformation
\[
  w \;=\; e^{s}, \qquad s=\log r+i\theta ,
\]
which we call the \emph{hidden-number coordinate system}. The map sends the
$s$-plane to the punctured $w$-plane, but it does more than rename variables:
the imaginary direction becomes an angle recorded modulo $2\pi$, so the plane
unrolls into a universal cover whose sheets are indexed by integers $k$. In
these coordinates the Dirichlet series of $\zeta$ becomes a sum of vectors of
locked length $(n+1)^{-1/2}$ rotating with angular velocity $\log(n+1)$, and the
functional equation becomes a pure inversion $w\mapsto e/w$. Geometry that is
invisible in the $s$-plane becomes visible---and provable---in the $w$-plane.

Every theorem stated below carries the name of its Lean formalization
(typeset as \Lean{name}); each of them compiles in the repository without
\Lean{sorry}. Our contribution is deliberately two-fold, and we insist on the
distinction throughout:

\begin{enumerate}
\item \textbf{Structural theorems.} A rigorous corpus concerning the geometry of
the coordinate system itself: the identity of the critical circle and of the
number $1/2$ (Section~\ref{sec:criticalcircle}); the reflection--inversion
correspondence and a closed equivalence cycle (Section~\ref{sec:reflection});
zero-location theory (Section~\ref{sec:zerolocation}); reflection pairs and
height interfaces (Section~\ref{sec:pairs}); the phase mechanism
(Section~\ref{sec:phase}); argument-principle atoms and single-zero certificates
(Section~\ref{sec:winding}); growth decomposition and conditional existence
chains (Section~\ref{sec:growth}); and energy divergence along the Hardy
direction (Section~\ref{sec:energy}).
\item \textbf{A proof roadmap with explicit gaps.} The coordinate framework
does \emph{not} prove the Riemann hypothesis. What it does is compress all
remaining difficulty into explicitly named analytic inputs---most notably a
height-supply interface (Definition~\ref{def:supply}) equivalent in strength to
$N(T)\to\infty$---and organize everything else into automatically closing
bridges. The honest boundaries are catalogued in Section~\ref{sec:honesty}.
\end{enumerate}

We emphasize already in the introduction, and repeat in the conclusion: the
phrase ``solution in hidden-number coordinates'' in the title refers to the
solution \emph{framework}---the coordinate system, the formalized structural
theory, and the roadmap. The final proof of the Riemann hypothesis still depends
on analytic inputs that are isolated explicitly in this paper but not supplied.

All formalized results were developed in the interactive theorem prover Lean
\citep{deMoura2015} against its mathematical library \citep{MathlibCommunity2020}.

\section{The hidden-number coordinate system}\label{sec:coords}

\begin{definition}[hidden-number coordinates]\label{def:coords}
The map $\Phi(s)=e^{s}$ sends $s=u+i\theta$ to $w=re^{i\varphi}$ with
$r=e^{u}$ and $\varphi=\theta \bmod 2\pi$. Because the angle forgets the sheet,
the honest domain of $\Phi$ is the universal cover of $\mathbb{C}\setminus\{0\}$;
a point of the cover is a pair $(w,k)$ with $k\in\mathbb{Z}$ recording the
winding number, and $\theta=2\pi k+\arg w$.
\end{definition}

The transport of $\zeta$ to these coordinates is given by the principal-branch
formula.

\begin{proposition}[principal branch;\,\Lean{expZeta_phase_principal}]\label{prop:principal}
On the punctured $w$-plane, define $\widehat{\zeta}(re^{i\varphi})
:=\zeta(\log r+i\varphi)$ with the real logarithm. Then $\widehat{\zeta}$ agrees
with $\zeta$ under the coordinate change wherever the right-hand side is
defined, and it constitutes the principal sheet of the lifted function.
\end{proposition}

The sheets interact by exact periodicity, and particular sheets fold onto rays.

\begin{proposition}[multi-sheet periodicity and folding;\,
\Lean{phase_periodic}, \Lean{phase_fold_neg_axis},
\Lean{phase_fold_many_sheets}]\label{prop:periodic}
For every integer $k$ and every admissible point,
\[
  \widehat{\zeta}_{k}(re^{i\varphi})
  \;=\; \zeta\bigl(\log r+i(\varphi+2\pi k)\bigr)
  \;=\; \widehat{\zeta}_{0}(re^{i\varphi}),
\]
so consecutive leaves differ only by the deck transformation $\theta\mapsto
\theta+2\pi$. The leaf starting at $\theta=\pi$ (more generally
$\pi+2\pi k$) folds onto the negative real axis of the $w$-plane.
\end{proposition}

Each turn of the spiral carries a genuine copy of the phase data.

\begin{theorem}[branch over every turn;\,\Lean{envelope_universal_cover_branch},
\Lean{helix_continuation}]\label{thm:branches}
For every $w\neq 0$ and every integer $k$ there exists a principal-branch phase
preimage on the $k$-th leaf; consequently the helical continuation of the
lifted function extends through every turn of the universal cover.
\end{theorem}

\begin{remark}
The multivalued-logarithm packaging itself,
\Lean{EnvelopeUniversalCoverMultiLog}, is currently maintained in the repository
as a declared structure rather than a proved theorem; the isomorphism statement
below deliberately avoids depending on it.
\end{remark}

The structural soundness of the whole picture is summarized by an isomorphism
of descriptions.

\begin{theorem}[cover--plane correspondence;\,\Lean{zetaCoverIsomorphism}]
\label{thm:coveriso}
The following data constitute a structure isomorphism between the covering
description and the classical complex-plane description: a bijection
\Lean{covers} between sheets and turns; an injection \Lean{radial} on radii;
value agreement \Lean{zeta_eq}; spiral continuation \Lean{branch}; and
multi-leaf folding \Lean{periodic}. Every constituent field is a proved
theorem; the assembly contains no admitted steps.
\end{theorem}

\begin{remark}
A useful slogan, kept deliberately crude: the stage (the universal cover with
its leaves and folding) is complete; whether actors (zeros) stand on it is a
separate question answered nowhere by the stage alone.
\end{remark}

\section{The critical circle and the identity of \texorpdfstring{$1/2$}{1/2}}
\label{sec:criticalcircle}

\begin{theorem}[critical line maps to a circle;\,\Lean{criticalLine_circle}]
\label{thm:critcircle}
For every real $t$,
\[
  \Norm{e^{\tfrac12+it}} \;=\; \sqrt{e}\, .
\]
Hence the image of the critical line $\Re s=1/2$ under the hidden-number map is
exactly the circle of radius $\sqrt{e}$ centered at the origin.
\end{theorem}

The radius itself is not a convention but a logarithmic identity.

\begin{proposition}[\Lean{log_sqrt_exp_one}]\label{prop:logsqrt}
$\log(\sqrt{e}) = \tfrac12$.
\end{proposition}

The decisive structural fact is that the critical circle is not merely
\emph{a} distinguished radius: among all concentric circles it is the unique
one preserved by the inversion induced by the functional equation. Write the
inversion as $\iota(w)=e/w$.

\begin{theorem}[unique fixed circle;\,\Lean{envelope_inversion_fixed_circle}]
\label{thm:fixedcircle}
A circle $\abs{w}=\rho$ with $\rho>0$ satisfies
$\abs{\iota(w)}=\abs{w}$ on the circle if and only if $\rho^{2}=e$. The
equation $\rho^{2}=e$ has a unique positive real solution, namely
$\rho=\sqrt{e}$.
\end{theorem}

\begin{corollary}[why $1/2$ is structural;\,from
Theorems~\ref{thm:fixedcircle} and \ref{thm:critcircle},
Proposition~\ref{prop:logsqrt}]\label{cor:onesecond}
If the reflection symmetry of the functional equation is to fix a scale---that
is, to preserve some circle centered at the origin---then that circle is
uniquely determined to be $\abs{w}=\sqrt{e}$, and its pullback to the $s$-plane
is uniquely the line $\Re s=1/2$. Within the hidden-number coordinates, the
location $1/2$ is thereby upgraded from a definition to a consequence of
reflection geometry: it is the only candidate position, not one among several.
\end{corollary}

On the fixed circle, the inversion acts as nothing other than conjugation.

\begin{theorem}[inversion equals conjugation on the circle;\,
\Lean{envelope_inversion_eq_conj_on_circle}]\label{thm:invconj}
For every $w$ with $\abs{w}=\sqrt{e}$,
\[
  \frac{e}{w} \;=\; \overline{w}\, .
\]
This is an identity holding at \emph{all} points of the circle; it is not a
statement about zeros.
\end{theorem}

\begin{remark}[two layers of ``position'']\label{rem:twolayers}
Corollary~\ref{cor:onesecond} proves \emph{layer 1}: the \emph{identity} of the
critical location---if reflection symmetry fixes a circle, it must be this one.
It does \emph{not} prove \emph{layer 2}: that any zero lies on that circle.
Combining layers, Theorem~\ref{thm:invconj} yields only the conditional
statement ``\emph{if} a zero $w$ lies on the critical circle, then its
functional-equation partner $e/w$ coincides with $\overline{w}$, i.e.\ zeros on
the line are naturally self-paired''. Whether any point lies there is the
content of the mechanism hypothesis discussed in
Section~\ref{sec:reflection}, and it is not delivered by the constant $e$
alone. Likewise, the circle is a continuum with uncountably many points---an
axiom of geometry---while zeta's nontrivial zeros form a discrete set; no
cardinality information passes from the former to the latter.
\end{remark}

\section{Reflection--inversion correspondence}\label{sec:reflection}

The functional equation pairs $s$ with $1-s$. In hidden-number coordinates this
becomes an inversion about the critical circle.

\begin{theorem}[reflection as inversion;\,\Lean{envelope_inversion_map}]
\label{thm:map}
Under $w=e^{s}$, the reflection $s\mapsto 1-s$ corresponds exactly to the map
\[
  w \;\longmapsto\; \frac{e}{w}\, :
  \qquad e^{\,1-s} \;=\; \frac{e}{e^{s}}\, .
\]
Moreover $s\mapsto 1-s$ is an involution
(\Lean{reflection_involutive}), and so is $w\mapsto e/w$.
\end{theorem}

With the dictionary of Theorem~\ref{thm:map} the classical equivalence web
around the hypothesis organizes into a closed cycle of four statements:
\begin{itemize}
\item[\textup{(A)}] no zeros outside the far region ($\Re s\geq 1$);
\item[\textup{(B)}] all nontrivial zeros lie on the critical circle
      ($\abs{w}=\sqrt{e}$);
\item[\textup{(C)}] the envelope radius is locked: $\abs{e^{s}}=\sqrt{e}$ at
      zeros;
\item[\textup{(D)}] the classical Riemann hypothesis.
\end{itemize}

\begin{theorem}[closed mechanism cycle;\,\Lean{mechanism_cycle_closed}]
\label{thm:cycle}
Under the explicit mechanism assumption $h_0:\ 0<\Re s$ (zeros lie in the right
half-plane), the implications $A\Rightarrow B\Rightarrow C\Rightarrow D\Rightarrow
A$ are all proved, in each of the four directions.
\end{theorem}

\begin{remark}[what the cycle does and does not say]\label{rem:cycle}
Theorem~\ref{thm:cycle} is an \emph{equivalence canon}: it certifies that the
four formulations are mutually interchangeable, so that any future proof of one
instantly yields the others. It is \emph{not} a truth certificate for any of
them. Two caveats are recorded verbatim from the formalization. First, the leg
$A\Rightarrow B$ is proved outright, but the closing leg $D\Rightarrow A$ routes
through MathLib's declared object \Lean{RiemannHypothesis}; the closure
therefore leans on the hypothesis itself. Second, the input $h_0$ is a
classical fact not yet realized at the formal level. The cycle tells us the
statements agree; it does not tell us any of them holds.
\end{remark}

\section{Zero location theory}\label{sec:zerolocation}

The first two theorems exclude the two flanks of the critical strip.

\begin{theorem}[far side;\,\Lean{zero_free_outside_envelope}]\label{thm:farside}
If $\Re s\geq 1$ then $\zeta(s)\neq 0$. The proof transports the positivity of
the Euler product: each local factor is nonzero, hence so is the product
(cf.\ MathLib's \Lean{riemannZeta_ne_zero_of_one_le_re}).
\end{theorem}

\begin{theorem}[negative side;\,\Lean{zero_free_negative_side}]
\label{thm:negside}
If $\Re s\leq 0$ and $s\notin\{0,-2,-4,-6,\dots\}$, then $\zeta(s)\neq 0$; the
argument reflects the far-side result through the functional equation
(Theorem~\ref{thm:map}).
\end{theorem}

\begin{corollary}[critical strip;\,\Lean{nontrivial_zero_in_critical_strip}]
\label{cor:strip}
Every nontrivial zero satisfies $0<\Re s<1$.
\end{corollary}

Transported to the envelope, the strip becomes an annulus whose geometric mean
is the critical circle.

\begin{theorem}[annulus localization;\,
\Lean{nontrivial_zero_in_envelope_annulus}]\label{thm:annulus}
If $s$ is a nontrivial zero, then
\[
  1 \;<\; \Norm{e^{s}} \;<\; e\, ,
  \qquad\text{and}\qquad
  \Norm{e^{s}} = e^{\Re s}\, .
\]
The critical leaf $\Norm{w}=\sqrt{e}$ sits exactly at the geometric mean of the
annulus $1<\abs{w}<e$: the multiplicative midpoint between the inner and outer
radii is $\sqrt{1\cdot e}=\sqrt{e}$.
\end{theorem}

Trivial zeros, by contrast, are abundant and completely explicit.

\begin{theorem}[infinitely many trivial zeros;\,
\Lean{infinitely_many_trivial_zeros}]\label{thm:trivial}
The set of trivial zeros is infinite: the map $n\mapsto -2(n+1)$ is injective
and each value is a zero, $\zeta(-2(n+1))=0$. No analytic growth input of any
kind is used; the infinitude comes from an explicit formula.
\end{theorem}

\begin{remark}
The contrast between Theorems~\ref{thm:trivial} and \ref{thm:annulus} frames the
central asymmetry of the subject. Trivial zeros are infinite because we possess
a formula producing them; nontrivial zeros are located---confined to the
annulus---but no existence or infinitude follows from location alone. This
asymmetry drives the design of the interfaces in
Sections~\ref{sec:pairs} and \ref{sec:growth}.
\end{remark}

\section{Reflection pairs in the annulus}\label{sec:pairs}

\begin{proposition}[\Lean{inversion_preserves_annulus}]\label{prop:preserves}
The inversion $w\mapsto e/w$ preserves the annulus $1<\abs{w}<e$.
\end{proposition}

\begin{definition}[annular nontrivial zero;\,
\Lean{AnnularNontrivialZero}]\label{def:annzero}
An \emph{annular nontrivial zero} is the packaged object consisting of a point
$s$ lying in the critical strip together with the witness that
$\Norm{e^{s}}$ lies in the annulus. The constructor
\Lean{annularNontrivialZero_of_zero} promotes any nontrivial zero to this
object.
\end{definition}

\begin{proposition}[closure under reflection;\,
\Lean{annularNontrivialZero_reflects},
\Lean{annularNontrivialZero_iff_reflects},
\Lean{on_circle_reflects_on_circle}]\label{prop:closes}
The set of annular nontrivial zeros is closed under the reflection $s\mapsto
1-s$, and the closure is reversible (an ``iff''). On the critical circle the
reflection preserves the circle as a \emph{set}; this is set-preservation, not
a pointwise fixing of every point. Trivial negative-even zeros are excluded by
hypothesis from being mis-reflecting into the picture.
\end{proposition}

\begin{theorem}[existence of reflection pairs;\,
\Lean{exists_annular_nontrivial_zero_pair}]\label{thm:pairs}
Given one annular nontrivial zero, its reflection partner can be constructed,
producing a genuine pair $\{s,\,1-s\}$. The theorem is deliberately
\emph{conditional}: from one zero it produces two, never from none it produces
one. It does not imply infinitude.
\end{theorem}

To discuss heights one lifts finite rectangles into the cover.

\begin{definition}[strip windows;\,\Lean{hiddenStripWindow},
\Lean{hiddenLiftRectangle}]\label{def:window}
For $T>0$ the \emph{hidden strip window} is the lifted rectangle
$\{s:\ 0<\Re s<1,\ \abs{\Im s}<T\}$ inside the universal cover; it retains the
spiral angle $\theta=\Im s$ without folding. Its symmetric variant
\Lean{hiddenLiftRectangle} keeps the same leaf coordinates on all four sides.
\end{definition}

\begin{proposition}[boundary integrals;\,
\Lean{hiddenRectangleBoundaryIntegral},
\Lean{hiddenRectangleBoundaryIntegral_eq_zero_of_differentiableOn}]
\label{prop:boundaryint}
The four sides of the lifted rectangle assemble into a boundary integral in
hidden-number coordinates, and for any function holomorphic on the rectangle
this integral vanishes. The proof invokes the Cauchy--Goursat rectangle atom of
MathLib; nothing about $\zeta$ beyond holomorphicity is used.
\end{proposition}

Heights are then extracted and converted into infinitude---given the height
growth.

\begin{definition}[height interface;\,
\Lean{hiddenNontrivialZeroHeights},
\Lean{HiddenZeroHeightGrowth}]\label{def:heightgrowth}
Annular nontrivial zeros are mapped to their cover heights
$\abs{\Im s}$. \emph{Hidden zero height growth} is the assertion that this set
of heights is unbounded.
\end{definition}

\begin{theorem}[height growth yields infinitude;\,
\Lean{infinite_annular_nontrivial_zeros_of_height_growth},
\Lean{infinite_nontrivial_zeros_of_hidden_height_growth}]
\label{thm:heightgrowth}
If the hidden zero height grows, then the set of annular nontrivial zeros is
infinite, and consequently the set of nontrivial zeros of $\zeta$ in the
classical critical strip is infinite. Both implications are purely
set-theoretic bridges: bounded heights would confine zeros to a compact
rectangle admitting only finitely many, by discreteness; the proofs consume no
analysis of $\zeta$ whatsoever.
\end{theorem}

\begin{remark}
At the time of writing \Lean{HiddenZeroHeightGrowth} itself is \emph{not}
proved; supplying it is identified in Section~\ref{sec:growth} as the next
analytic input. What Sections 5--6 establish jointly is the full machinery
surrounding the black box: everything that can be said about zeros from
symmetry, topology and set theory alone has been said, and said formally.
\end{remark}

\section{The phase mechanism}\label{sec:phase}

On the critical leaf the Dirichlet series of $\zeta$ becomes a superposition of
rotating vectors of locked lengths.

\begin{proposition}[length lock;\,\Lean{critical_leaf_norm_locked}]
\label{prop:lock}
On the critical leaf ($r=\sqrt{e}$), every term of the series has modulus
exactly
\[
  (n+1)^{-1/2}\, .
\]
The length dimension is frozen; only phases rotate.
\end{proposition}

\begin{proposition}[absolute convergence divide;\,
\Lean{absolute_convergence_outside_critical_leaf},
\Lean{critical_leaf_not_absolutely_convergent}]\label{prop:convdiv}
Absolute convergence holds strictly outside radius $e$ (leaves farther out),
while the critical leaf $r=\sqrt{e}<e$ fails absolute convergence: the critical
leaf sits on the non-absolutely-convergent side of the divide.
\end{proposition}

\begin{proposition}[winding counts;\,\Lean{term_winding_number}]
\label{prop:windingcount}
The number of turns made by the $n$-th term over the natural phase parameter is
$\log(n+1)$---angular speed multiplied by leaf count---recorded as
\Lean{windingCount}.
\end{proposition}

Phase cancellation admits an exact combinatorial characterization.

\begin{definition}[finite phase cancellation;\,
\Lean{phaseCancelRelation}]\label{def:cancel}
A finite integer combination $\sum_i a_i\log(n_i+1)$ \emph{cancels} if it lies
in $2\pi\mathbb{Z}$: finitely many frequencies, weighted by integers, return to
their starting angle on the unit circle.
\end{definition}

\begin{proposition}[\Lean{phaseCancel_of_exp_one}]\label{prop:expone}
Cancellation implies $\exp\!\bigl(i\sum_i a_i\log(n_i+1)\bigr)=1$, i.e.\ the
product of unit vectors returns to $1$.
\end{proposition}

\begin{theorem}[cancellation equals integer power identity;\,
\Lean{phaseCancel_zero_iff_powProduct_one}]\label{thm:cancel}
Finite phase cancellation is equivalent to the integer multiplicative identity
\[
  \sum_i a_i\log(n_i+1)\in 2\pi\mathbb{Z}
  \quad\Longleftrightarrow\quad
  \prod_i (n_i+1)^{a_i} \;=\; 1\, .
\]
The characterization is exact, purely algebraic, and independent of
transcendence questions.
\end{theorem}

\begin{remark}[degeneracy of finite cancellation]\label{remark:degen}
Since $\sum_i a_i\log(n_i+1)=\log\prod_i (n_i+1)^{a_i}=\log q$ with $q$
rational, membership in $2\pi\mathbb{Z}$ forces the integer part to vanish
($e^{2\pi k}$ is rational only for $k=0$, an instance of
Lindemann--Weierstrass; the analytic half \Lean{exp_polynomial_approx} is
available in MathLib). Hence every cancelling combination is degenerate:
the shifted-window phase $-t\cdot\textup{(cancel-sum)}$
(\Lean{phaseShift_cancel_of_int}, \Lean{phaseShift_aligned_at_integer}) stays
identically $0$ along integer shifts, and alignment windows never climb to
unbounded heights by this route. Finite phase combinations cannot manufacture
abundance; this coordinate-internal negative result rules out one seductive
shortcut.
\end{remark}

A parallel mechanism runs through the alternating function $\eta$, valid on the
convergent side.

\begin{proposition}[alternating bridge;\,
\Lean{eta_eq_mul_zeta}, \Lean{eta_term_rotating_vector},
\Lean{eta_phase_alignment_condition}]\label{prop:eta}
For $\Re s>1$,
\[
  \eta(s) \;=\; \bigl(1-2^{1-s}\bigr)\,\zeta(s),
\]
and the $n$-th term of $\eta$ factors as $(-1)^{n}$ times a rotating vector.
Consequently, for $\Re s>1$: $\eta(s)=0$ if and only if the alternating sum of
rotating vectors vanishes.
\end{proposition}

\begin{definition}[\Lean{alignmentZero}]\label{def:align}
$s$ is an \emph{alignment zero} if its alternating rotating-vector sum (the
$\eta$ shape) equals zero.
\end{definition}

\begin{theorem}[\Lean{alignmentZero_iff_zeta_zero}]\label{thm:align}
For $\Re s>1$, alignment zero at $s$ is equivalent to $\zeta(s)=0$; the
rejection factor $1-2^{1-s}$ is nonzero there. Alignment zeros project onto
classical zeros.
\end{theorem}

\begin{remark}
Within the critical strip $0<\Re s<1$ the factor $1-2^{1-s}$ remains nonzero
(it vanishes only at real part $1$), and $\zeta=0\Longleftrightarrow\eta=0$
holds classically; but the formal transfer requires the conditionally convergent
$\eta$ series and Abel/function-equation regularization, tooling not yet
available in MathLib (see Section~\ref{sec:honesty}, item (v)). Finally, the
``phase budget'' on the critical circle is a genuinely operational quantity:
the per-term phase velocity has norm $(n+1)^{-1/2}\log(n+1)$
(\Lean{critical_leaf_phase_velocity_term_norm}), budgets accumulate
monotonically in $N$ (\Lean{critical_leaf_phase_velocity_budget_increasing})
with positive increments (\Lean{critical_leaf_phase_velocity_positive}). Yet
budget accumulation does not force a zero: the tail
$\sum_{n>N}(n+1)^{-1/2}e^{-i\theta\log(n+1)}$ does not tend to $0$ on the
critical circle, and controlling it is exactly the growth-input gap of
Section~\ref{sec:growth}.
\end{remark}

\section{Argument principle and winding atoms}\label{sec:winding}

\begin{theorem}[single-zero atom;\,
\Lean{argumentPrinciple_single_zero}]\label{thm:atom}
If the circle encloses the single point $w_{0}$ and excludes the origin's
singularity, then
\[
  \oint \frac{dz}{z-w_{0}} \;=\; 2\pi i\, .
\]
\end{theorem}

\begin{proposition}[circle parametrization;\,
\Lean{envelope_circle_param}]\label{prop:param}
The circle $\abs{z-c}=r$ in envelope-phase coordinates is parametrized by
$\theta\mapsto c+r e^{i\theta}$.
\end{proposition}

\begin{theorem}[phase winding counts zeros;\,
\Lean{phase_winding_equals_zero_count}]\label{thm:phasewind}
In logarithmic-derivative form, the winding of the phase along a contour equals
(the number of enclosed zeros) times $2\pi i$: one full turn of phase per zero.
\end{theorem}

\begin{proposition}[\Lean{zeroCount_normalized}]\label{prop:norm}
Normalizing removes the constant: $(2\pi i)^{-1}\oint dz/(z-w_{0})=1$.
\end{proposition}

\begin{theorem}[power-factor winding;\,
\Lean{powFactor_logDeriv_eq}, \Lean{winding_of_pow_factorization}]
\label{thm:powwind}
If $f(z)=(z-w_{0})^{m}g(z)$ with $g$ holomorphic and nonvanishing on the disk,
then
\[
  \oint \frac{f'}{f} \;=\; 2\pi i\, m\, .
\]
The proof composes logarithmic-derivative multiplicativity with power and
composition lemmas, circumventing delicate natural-subtraction power algebra.
\end{theorem}

\begin{remark}[completeness of the argument principle]
The full argument principle is Theorem~\ref{thm:powwind} plus enumeration of a
finite zero set with multiplicity summation; neither the enumeration nor the
Riemann--von Mangoldt counting function is formalized, the latter requiring
Stirling-type growth estimates (Section~\ref{sec:growth}).
\end{remark}

The winding atoms become useful through certificates that package a zero with
everything needed downstream.

\begin{definition}[simple zero certificate;\,
\Lean{ZetaSimpleZeroCertificate}]\label{def:cert}
A \emph{simple zero certificate} consists of an open set $U$, a point $w\in U$,
and a function $g$ holomorphic and nonzero on $U$ such that
\[
  \zeta(z) \;=\; (z-w)\, g(z) \quad\text{for } z\in U\, .
\]
Both the zero $w$ and the factor $g$ are explicit parameters, so that every
field of the structure is a proposition.
\end{definition}

\begin{theorem}[certificate consequences;\,
\Lean{certificate_gives_zero},
\Lean{zeta_winding_eq_two_pi_i_of_certificate},
\Lean{annulusZeroWitness_of_zetaCertificate},
\Lean{exists_annular_pair_of_zetaCertificate}]\label{thm:certchain}
From any simple zero certificate the following follow automatically:
(i) $\zeta(w)=0$ (evaluate the factorization at $w$);
(ii) the logarithmic derivative winds once,
$\oint\zeta'/\zeta=2\pi i$ over suitable contours in $U$;
(iii) the certificate induces an annulus zero witness in the sense of
Section~\ref{sec:growth};
(iv) hence an annular reflection pair exists.
All downstream connections close with no further input: whoever supplies one
certificate inherits the entire chain for free.
\end{theorem}

\section{Growth decomposition and existence interfaces}\label{sec:growth}

\begin{theorem}[growth decomposition;\,
\Lean{zeta_growth_decomposition}]\label{thm:decomp}
Coordinate-independent and already proved in MathLib:
\[
  \zeta(s) \;=\; \frac{\Lambda_{0}(s)}{\Gamma_{\mathbb{R}}(s)}\, ,
\]
where $\Lambda_{0}$ denotes the completed zeta function
(\Lean{completedRiemannZeta}) and $\Gamma_{\mathbb{R}}$ the real Gamma factor.
\end{theorem}

The norms of the factors admit exact readings in hidden-number coordinates.

\begin{proposition}[\Lean{gammaR_norm_decomposition},
\Lean{gammaR_norm_in_envelope_coords},
\Lean{zeta_norm_in_envelope_coords}]\label{prop:normread}
Purely algebraically,
\[
  \Norm{\Gamma_{\mathbb{R}}(s)}
  \;=\; \pi^{-\Re s/2}\,\Norm{\Gamma(s/2)}\, ;
\]
on the critical leaf $\abs{w}=\sqrt{e}$ this specializes to
$\pi^{-1/4}\Norm{\Gamma((\log\sqrt{e}+i\theta)/2)}$, and consequently
\[
  \Norm{\zeta(\log r+i\theta)}
  \;=\;
  \frac{\Norm{\Lambda_{0}(\log r+i\theta)}}
       {\pi^{-\log r/2}\,\Norm{\Gamma((\log r+i\theta)/2)}}\, .
\]
These identities mark the precise point where growth enters the coordinate
picture; they predict, but do not prove, boundedness of the numerator.
\end{proposition}

Local analysis converts integrals into witnesses.

\begin{proposition}[\Lean{analyticAt_riemannZeta}]\label{prop:analytic}
$\zeta$ is analytic at every $s\neq 1$, proved via excised-neighborhood
differentiability glued by continuity.
\end{proposition}

\begin{theorem}[\Lean{logDeriv_circle_integral_eq_zero_of_zero_free}]
\label{thm:zerofreeint}
If $\zeta$ has neither zero nor pole in the closed disk, then the logarithmic-
derivative integral over the bounding circle vanishes.
\end{theorem}

\begin{theorem}[contrapositive existence bridge;\,
\Lean{exists_closedBall_zero_of_logDeriv_integral_ne_zero}]
\label{thm:existbridge}
If the logarithmic-derivative circle integral is nonzero, then $\zeta$ has a
zero in the closed disk. This is the exact logical contrapositive of
Theorem~\ref{thm:zerofreeint}; the formalization does \emph{not} claim that any
such nonzero integral occurs.
\end{theorem}

\begin{definition}[\Lean{AnnulusZeroWitness},
\Lean{CriticalStripWitness}]\label{def:witness}
A \emph{finite-disk witness} packages a disk with a nonzero winding; the
critical-strip variant adds the constraint that the disk lies inside the strip.
These are verifiable objects, and the code contains no claim that one exists.
\end{definition}

Everything above feeds a single, sharply isolated black box.

\begin{definition}[height supply;\,
\Lean{ZeroHeightSupply}]\label{def:supply}
\emph{Height supply} is the assertion: above every height $H$ there exists a
critical-strip zero. It is the direct weakening of $N(T)\to\infty$; it is
formalized as an explicit definition, neither an axiom nor a proved theorem.
\end{definition}

\begin{theorem}[supply chain;\,
\Lean{hiddenZeroHeightGrowth_of_zeroHeightSupply},
\Lean{infinite_nontrivial_zeros_of_zeroHeightSupply}]\label{thm:chain}
Height supply implies hidden zero height growth
(Definition~\ref{def:heightgrowth}), which implies infinitely many nontrivial
zeros. The chain is unconditional in its logic and conditional in its premise.
\end{theorem}

Certificates refine supply into a locally checkable currency.

\begin{definition}[unbounded certificate family;\,
\Lean{UnboundedZetaCertificateAssumption}]\label{deffamily}
The \emph{unbounded certificate family} is the conditional predicate
(formalized as a definition, not an axiom): above every height there exists a
simple-zero certificate instance (Definition~\ref{def:cert}) with disk inside
the critical strip.
\end{definition}

\begin{theorem}[family bridges;\,
\Lean{zeroHeightSupply_of_certificateFamily},
\Lean{classicalZeroCountGrowth_of_zeroHeightSupply},
\Lean{ModulatedZetaFamily},
\Lean{zeroHeightSupply_of_modulatedFamily},
\Lean{infinite_nontrivial_zeros_of_modulatedFamily}]\label{thm:family}
The following implications hold:
\[
  \textup{certificate family}
  \Longrightarrow \textup{height supply}
  \Longrightarrow \textup{height growth}
  \Longrightarrow \infty\text{ many nontrivial zeros},
\]
and, in parallel,
\[
  \textup{height supply}
  \Longrightarrow \textup{classical count growth}
  \Longrightarrow \textup{hidden count growth}
  \Longrightarrow \infty\text{ many nontrivial zeros}.
\]
The \emph{modulated elementary family}
(\Lean{ModulatedZetaFamily}) repackages the same interface: one certificate per
height, ``the same elementary object modulated without bound'', and inherits
the full chain by reuse.
\end{theorem}

\begin{remark}[the honest reading of Section 9]\label{remark:family}
What has been proved is every arrow. What has \emph{not} been proved is the
existence of a single certificate instance, or of height supply, or of
$N(T)\to\infty$. The unique missing analytic input has been compressed,
isolated and named---but it remains external. Supplying it is precisely the
Stirling/completed-zeta program recorded in Section~\ref{sec:honesty},
item (ii).
\end{remark}

\section{Energy divergence and the Hardy direction}\label{sec:energy}

\begin{proposition}[\Lean{critical_leaf_diagonal_energy_term}]\label{prop:diagterm}
On the critical leaf, squaring the locked lengths
(Proposition~\ref{prop:lock}) gives the diagonal energy of the $n$-th term:
\[
  \Norm{a_n}^{2} \;=\; (n+1)^{-1}\, .
\]
\end{proposition}

\begin{theorem}[\Lean{critical_leaf_diagonal_energy_diverges}]
\label{thm:diverge}
The corresponding series $\sum_n (n+1)^{-1}$ is the harmonic series and does
not converge: the diagonal energy on the critical leaf diverges. This is the
seed of abundance in the spirit of \citet{Hardy1914}---obtained with no appeal
to Stirling asymptotics.
\end{theorem}

Making the energy computation quantitative requires the rotating-integral
template.

\begin{definition}[\Lean{rotVecOnLine}, \Lean{dirichletPartialSumLine}]
\label{def:template}
The rotating vector on the critical line, $v_n(t)=(n+1)^{-1/2}
e^{-it\log(n+1)}$, and the associated finite Dirichlet partial sums are defined
explicitly as the technical template for mean-value computations.
\end{definition}

\begin{proposition}[rotating integral atom;\,
\Lean{integral_rotating_atom}]\label{prop:intatom}
For $\Delta\neq 0$,
\[
  \int_{-T}^{T} \exp(i\Delta x)\,dx
  \;=\; \frac{e^{i\Delta T}-e^{-i\Delta T}}{i\Delta}\, .
\]
The antiderivative is certified directly via \Lean{HasDerivAt}; the
corresponding MathLib lemma is temporarily unreachable during the ongoing
module-system migration.
\end{proposition}

\begin{theorem}[single-term energy integral;\,
\Lean{energy_single_term_integral}]\label{thm:energyint}
Over the height window of length $2T$,
\[
  \int \Norm{v_n}^{2} \;=\; 2T\,(n+1)^{-2\sigma}\, .
\]
This is the exact place where ``logarithmic accumulation'' enters computation:
after passing to the diagonal sum, $\sum_{n}(n+1)^{-2\sigma}$ diverges at
$\sigma=1/2$ precisely as the harmonic series of
Theorem~\ref{thm:diverge}, and converges for $\sigma>1/2$.
\end{theorem}

\begin{remark}[the mean-formula gap]\label{remark:mean}
Classically, Hardy's theorem \citep{Hardy1914} concludes infinitely many zeros
on the critical line from the mean formula
$\int_{0}^{T}\abs{\zeta(1/2+it)}^{2}dt\asymp T\log T$. In the present framework
the diagonal part of that computation is formalized
(Theorem~\ref{thm:energyint}); what is missing is the cross-term control and
the limit exchange (Abel/uniform-convergence arguments), together with Hardy's
``finitely many zeros lead to contradiction'' step. MathLib possesses neither
the mean formula nor the surrounding machinery; the gap is recorded, not
papered over. Note also what energy divergence is and is not: it is fuel, not an
engine. Persistent quadratic energy of rotating terms constrains the size of
values on the critical line; it neither produces interference zeros by itself,
nor says anything about zeros off the line.
\end{remark}

Two further elementary inputs sharpen the picture on the Hadamard side.

\begin{theorem}[recursion lower bound;\,
\Lean{Real.gamma_ge_pow_mul_self}]\label{thm:gamma}
For all $x\geq 1$ and natural $n$,
\[
  \Gamma(x+n) \;\geq\; x^{n}\,\Gamma(x)\, ,
\]
proved by induction iterating $\Gamma(x+1)=x\,\Gamma(x)$. Completely
elementary; no Stirling estimate involved. This supplies the real-axis lower
bound of the Hadamard chain and demonstrates that, on that route, Stirling is
not a necessary passage.
\end{theorem}

\begin{remark}[elementary carriers and the Hadamard chain]
\label{remark:xientire}
The entire-function carrier of the nontrivial zeros has landed:
$\xi$ is realized as $\tfrac12 s(s-1)\Lambda_{0}(s)+\tfrac12$
(\Lean{xiEntire}), making entirety and symmetry pure polynomial consequences
via MathLib's pole-free auxiliary \Lean{completedRiemannZeta_0}, with the
explicit product formula $\xi=\tfrac12 s(s-1)\Gamma_{\mathbb{R}}\zeta$
(\Lean{xiEntire_eq_mul_zeta}); since $\Gamma_{\mathbb{R}}$ is nonzero on the
right half-plane, the zero bridge $\xi=0\Longleftrightarrow\zeta=0$ inside the
critical strip follows from this formula by inspection. The Hadamard combination
chain is mapped segmentwise: Borel--Carath\'eodory is available in MathLib
(\Lean{Complex.borelCaratheodory}), the $\Gamma$ lower bound is
Theorem~\ref{thm:gamma}, and the polynomial-forcing and real-axis contradiction
ends are light; the unformalized middle consists of the strip bound for $\zeta$
via quantified $\eta$/Abel estimates, coarse upper bounds for $\Gamma$ by
integral splitting, and their composition into the Borel--Carath\'eodory
application. Growth upper bounds would in principle need only
$\abs{\Gamma(s)}\leq e^{O(\abs{s}\log\abs{s})}$---again avoiding Stirling---but
until that chain closes, no infinitude claim is made from this direction.
Finally, the zero-isolation toolkit contains one proven exclusion tool
(\Lean{no_zero_of_boundary_dominance}: $\abs{h/g}\leq q<1$ on the circle forces
$g+h$ zero-free inside, via the maximum modulus principle; specialized to
$\zeta$ in \Lean{zeta_ne_zero_of_boundary_unit_bound}), while the existence
half---dominant linear term forcing exactly one zero---requires
winding-number integrality and homotopy invariance, absent from MathLib, so a
full Rouch\'e awaits infrastructure.
\end{remark}

\section{Honest boundaries and conclusion}\label{sec:honesty}

This section is the load-bearing wall of the paper: it enumerates what has
\emph{not} been proved, with the same precision as Sections 2--10 enumerate
what has.

\subsection*{(i) Position versus abundance: orthogonality}
The constant $e$ enters the critical structure at exactly three places, all
from one reflection/inversion mechanism: the fixed circle has radius
$\sqrt{e}$ (Theorem~\ref{thm:fixedcircle}); $\abs{w}=e$ is the outer annulus
boundary and the Dirichlet absolute-convergence divide
(Proposition~\ref{prop:convdiv}); and $\log\sqrt{e}=1/2$
(Proposition~\ref{prop:logsqrt}) selects the critical location. This explains
the \emph{positional identity} of $1/2$---a question the classical framework
cannot even pose. It explains \emph{nothing} about abundance: a circle is
continuous geometry carrying no cardinality information about the discrete
zero set upon it; the unique honest infinitude interfaces remain
Definitions~\ref{def:heightgrowth} and \ref{def:supply}. Any reasoning from
self-conjugacy of the circle to infinitely many zeros lacks a middle step.

\subsection*{(ii) The $N(T)\to\infty$ input is not supplied}
Unconditional infinitude of nontrivial zeros requires either the
Stirling-asymptotics-plus-completed-zeta-boundedness program behind
Riemann--von Mangoldt counting, or the Hardy mean formula
(Remark~\ref{remark:mean}). Neither is formalized; predictions such as
Stirling asymptotics and boundedness of the completed zeta remain open
boundaries explicitly marked as drafts in the repository. Until then
\Lean{ZeroHeightSupply} (Definition~\ref{def:supply}) is the single named
black box, and the project does not possess an unconditional proof that
infinitely many nontrivial zeros exist.

\subsection*{(iii) Recorded verdict: infinitely many circle points do not
give zeros}
The judgment is recorded formally in the repository: ``the circle has
infinitely many points'' (continuity axiom) does not imply ``$\zeta$ has a
single zero on the circle'', let alone infinitely many. Geometric continuity
supplies no zero instances; exclusion-and-location results
(Theorems~\ref{thm:farside}--\ref{thm:annulus}) likewise cannot manufacture
existence. The verdict is retained as a standing guard against a recurring
category error.

\subsection*{(iv) Three rejected shortcuts}
Three tempting jumps have been examined and refuted inside the coordinate
system itself. First: circle-continuity implying zero existence (rejected, item
(iii) above). Second: finite phase combinations implying unbounded heights---
refuted by the degeneracy analysis of Remark~\ref{remark:degen}, where every
cancelling combination collapses to a zero shift. Third: the self-conjugate
circle implying infinitely many zeros---same orthogonality failure as item (i).
(A related dead end is documented alongside: angle modulation cannot grow to
cover tails, since varying $t$ rotates existing terms but creates no new
frequencies or weights; frequency-domain growth reduces back to certificate
families.) The genuine constructive route---boundary-nonvanishing plus nonzero
contour winding inside finite-height sectors, promoted along an unbounded
sequence---is articulated as a candidate requiring the very supply input of
item (ii); the coordinate system contributes organization and verifiable
witnesses, not zeros from thin air.

\subsection*{(v) Missing MathLib infrastructure}
The following absences were reconnoitered concretely: interval arithmetic
(blocking strict numerical certification of Hardy-type sign changes; floating
point cannot certify Lean proofs---the workaround is
Definition~\ref{def:cert}, reducing numerical endpoints to pure predicates);
conditionally convergent $\eta$ series with Abel/function-equation
regularization (blocking the critical-strip transfer of
Theorem~\ref{thm:align}); winding-number integrality and homotopy invariance
(blocking the existence half of Rouch\'e, Remark~\ref{remark:xientire});
mean-square formulas for $\zeta$ on the line (blocking
Remark~\ref{remark:mean}); and transient module-system migrations hiding
individual integral lemmas (worked around by direct
\Lean{HasDerivAt} certification, Proposition~\ref{prop:intatom}).

\subsection*{Conclusion}
The hidden-number coordinate system $w=e^{s}$ converts the landscape of the
Riemann hypothesis into geometry: the critical line becomes the unique
inversion-fixed circle, the number $1/2$ becomes the forced logarithm of its
radius, the functional equation becomes an inversion acting as conjugation on
that circle, zeros become confined to an annulus centered geometrically on it,
and the analytic difficulties become explicit, named, and isolated interfaces
(Definitions~\ref{def:cert}, \ref{def:supply},
\ref{deffamily}). Every arrow between these objects is machine-checked; the
missing premises are stated rather than smuggled. The title's promise is thus
fulfilled in the precise sense this paper defends: a solution
\emph{framework}---coordinates, structural theorems, and a complete
roadmap---whose final step, the supply of the isolated analytic input, remains
open work. We believe that compressing a century-old problem into one
explicitly typed black box, surrounded on all sides by verified structure, is
itself a contribution: it tells future efforts exactly where to push.

\begin{acknowledgement}
The author thanks the Lean and MathLib communities for the library of formal
mathematics on which every theorem of this paper rests, and the maintainers of
the J-SFdS \LaTeX{} class.
\end{acknowledgement}

\bibliography{biblio}
\end{document}
